% !TEX program = pdflatex
\documentclass[11pt,a4paper]{article}

\usepackage[utf8]{inputenc}
\usepackage[T1]{fontenc}
\usepackage[french]{babel}
\usepackage{geometry}
\usepackage{array}
\usepackage{tabularx}
\usepackage{booktabs}
\usepackage{enumitem}
\usepackage{hyperref}
\usepackage{longtable}

\geometry{margin=2.2cm}
\setlength{\parindent}{0pt}
\setlength{\parskip}{0.5em}

\title{Cahier des charges\\
Avion de cartographie aérienne}
\author{Sofia Mabrouk, Emylie Hiron, Alain Piallat, Melvin Sayavong}
\date{2e année / 2026-2027}

\begin{document}

\maketitle
\tableofcontents
\newpage

\section{Déroulé du projet}

\textbf{Système embarqué de cartographie aérienne : avion RC + caméra stabilisée + acquisition et traitement des données}

Ce projet vise à appliquer les notions d'électronique et d'informatique des semestres précédents en développant 
un système embarqué capable d'acquérir des images du sol depuis un avion radiocommandé, de les associer à des 
informations de position, puis de produire une représentation cartographique de la zone survolée.

L'avion radiocommandé constitue principalement une plateforme expérimentale. Sa conception mécanique relève de 
l'aspect hobby du projet et ne constituent pas le coeur de l'ingénierie demandée.

Les deux objectifs principaux sont :

\begin{enumerate}
    \item \textbf{Commande de vol} : l'avion sera contrôlée par une télécommande.
    \item \textbf{Intégrer une charge utile électronique embarquée} : une caméra, un système GPS et un système de stockage.

    \item \textbf{Traiter les données de la caméra pour faire une orthomosaïque} à partir d'un logiciel déjà existant permettant de synchroniser les mesures, de reconstruire la trajectoire et d'assembler les photographies.
\end{enumerate}

Le projet sera développé de manière incrémentale. Une première version sera le test bench de tous les composants de 
l'avion et la charge utile, sans le corps de vol. L'attention sera concentrée sur le fonctionnement des moteurs, de 
la caméra et du système d'acquisition/exploitation. 
Une deuxième version intégrera la stabilisation de la caméra. La version finale devra permettre de produire automatiquement une représentation cartographique à partir des données acquises pendant un vol d'avion.

\section{Cahiers des charges}

\subsection{Cahier des charges du système électronique embarqué}

F: Fonction principale \\
M: Mécanique \\
E: Electronique \\
P: Perspectives, i.e. optionnelles, à faire si temps disponible et les autres charges respectées.

\begin{longtable}{|p{1cm}|p{3cm}|p{\dimexpr\textwidth-1cm-3cm-4\tabcolsep-4\arrayrulewidth\relax}|}
\hline
\textbf{ID} & \textbf{Catégorie} & \textbf{Exigence} \\ \hline
\endfirsthead

\hline
\textbf{ID} & \textbf{Catégorie} & \textbf{Exigence} \\ \hline
\endhead

F1 & Fonction principale &
Le système doit acquérir des images du sol pendant le vol et les enregistrer dans un format exploitable par le logiciel de cartographie. \\ \hline

F2 & Acquisition caméra &
La caméra doit permettre une cadence d'acquisition adaptée à la vitesse de déplacement de l'avion. \\ \hline

F3 & Stabilité &
L'acquisition s'arrête quand l'avion pivote. \\ \hline

F4 & Synchronisation &
Chaque image doit pouvoir être associée à une position (on aurait besoin de la position GPS, la latitude maybe, une vitesse aussi?). \\ \hline

F6 & Télécommande de l'avion &
L'avion sera radiocommandé par une télécommande préexistante. \\ \hline

F7 & Stockage &
Le système doit disposer d'une capacité de stockage suffisante pour une mission complète (SD card). \\ \hline

F8 & Exploitation des données &
Les données de la caméra pourront être traitées à l'aide d'un logiciel déjà existant pour créer une orthomosaïque. \\ \hline

M1 & Conception matérielle &
L'avion doit pouvoir voler et supporter la charge utile. \\ \hline

M2 & Intégration mécanique &
La charge utile doit pouvoir être fixée de manière fiable sur l'avion et protéger les composants électroniques contre les vibrations et les mouvements normaux du vol. \\ \hline

E1 & Gestion de l'énergie &
La consommation de la charge utile doit être mesurée et l'autonomie disponible doit être compatible avec une mission (durée d'un vol?). \\ \hline

E2 & Robustesse électrique &
Les alimentations doivent être filtrées afin de limiter les perturbations électriques liées au moteur, aux servomoteurs et aux autres équipements embarqués. \\ \hline

E3 & Système de caméra &
Un système  devra assurer la gestion de la caméra, du stockage, et potentiellement de la stabilisation de la caméra. \\ \hline

E4 & Commandes &
Un système gèrera la communication de la télécommande et de l'avion (pour les moteurs, ailes, train d'atterrissage). \\ \hline

P1 & Transmission des données &
Le système pourra transmettre en temps réel un retour caméra à l'utilisateur (faible débit).\\ \hline

P2 & Asservissement de l'avion &
Fonction 'stabilisation de l'avion' dans le quel il avance en ligne droite pendant x mètres/secondes \\ \hline

\end{longtable}

\begin{center}
Table 1 - Cahier des charges fonctionnel du système électronique embarqué
\end{center}

\end{document}
